\section{AHSME 1999}
        \begin{enumerate}
        	\item (\textit{AHSME 1999}) $1 - 2 + 3 -4 + \cdots - 98 + 99 =$
$\mathrm{(A) \ -50 } \qquad \mathrm{(B) \ -49 } \qquad \mathrm{(C) \ 0 } \qquad \mathrm{(D) \ 49 } \qquad \mathrm{(E) \ 50 }$

 \textbf{Answer: }E



\textbf{Solution 1}

If we group consecutive terms together, we get $(-1) + (-1) + \cdots + 99$, and since there are 49 pairs of terms the answer is $-49 + 99 = 50 \Rightarrow \mathrm{\textbf{(E)}}$.

 



\textbf{Solution 2}

( Similar to Solution 1 )If we rearranged the terms, we get $1+3-2+5-4 \cdots + 99-98$ then $1 + 1 + \cdots + 1$, and since there are 49 pairs of terms and the $1$ in the beginning the answer is $1+49 = 50 \Rightarrow \mathrm{\textbf(E)}$.

 



\textbf{Solution 3}

Let $1 - 2 + 3 -4 + \cdots - 98 + 99 = S$.

 Therefore, $S=99-98+97-\cdots -4+3-2+1$

 We add:

 $2S=100-100+100-100\cdots +100=100$
$S=50\Rightarrow \mathrm{\textbf{(E)}}$




\textbf{Solution 4}

We proceed with addition, 1 -2 + 3 -4.... Once done we find $\mathrm{\textbf{(E)}}$



	\item (\textit{AHSME 1999}) Which of the following statements is false?


 $\text{(A) All equilateral triangles are congruent to each other.}$
$\text{(B) All equilateral triangles are convex.}$
$\text{(C) All equilateral triangles are equiangular.}$
$\text{(D) All equilateral triangles are regular polygons.}$
$\text{(E) All equilateral triangles are similar to each other.}$

 \textbf{Answer: }A



\textbf{Solution}

An equilateral triangle is isosceles, and we find that $\angle A=\angle B=\angle C$ if we use the property of isosceles triangles that if two sides of a triangle are equal then the opposite angles are equal. Thus equilateral triangles are equiangular, and $C$ is true.

 Regular polygons are both equilateral and equiangular, and so are equilateral triangles are both equilateral (by definition) and equiangular (by the above argument). Thus equilateral triangles are regular polygons and $D$ is true. 

 Since all of the angles in an equilateral triangle are congruent, all equilateral triangles are similar by AAA similarity.  Thus, $E$ is true.

 Since $\angle A=\angle B=\angle C$ and $\angle A+\angle B+\angle C=180$, $\angle A=60^{\circ}$. Since no other angles are above $180^{\circ}$, all equilateral triangles are convex.  Thus, $B$ is true.

 This just leaves choice $\boxed{\mathrm{(A)}}$.  This is clearly false:  an equilateral triangle with side $1$ is not congruent to an equilateral triangle with side $2$.

 


	\item (\textit{AHSME 1999}) The number halfway between $1/8$ and $1/10$ is 


 $\mathrm{(A) \  } \frac 1{80} \qquad \mathrm{(B) \  } \frac 1{40} \qquad \mathrm{(C) \  } \frac 1{18} \qquad \mathrm{(D) \  } \frac 1{9} \qquad \mathrm{(E) \  } \frac 9{80}$

 \textbf{Answer: }E



\textbf{Solution 1}

To find the number halfway between $\frac{1}{8}$ and $\frac{1}{10}$, simply take the arithmetic mean, which is

 $\frac{\frac{1}{8}+\frac{1}{10}}{2}=\frac{\frac{9}{40}}{2}=\frac{9}{80}.$

 Thus the answer is choice $\boxed{E}.$




\textbf{Solution 2}

Note that $\frac{1}{10} = 0.1$ and $\frac{1}{8} = 0.125$.  Thus, the answer must be greater than $\frac{1}{10}$.  

 Answers $A$, $B$, and $C$ are all less than $\frac{1}{10}$, so they can be eliminated.  

 Answer $D$ is equivalent to $0.\overline{1}$, which is $0.0\overline{1}$ away from $0.1$, and is $0.013\overline{8}$ away from $0.125$.  These distances are not equal, eliminating $D$.

 Thus, $\boxed{E}$ must be the answer.  Computing $\frac{9}{80} = 0.1125$ as a check, we see that it is $0.1125 - 0.1 = 0.0125$ away from $\frac{1}{10}$, and similarly it is $0.125 - 0.1125 = 0.0125$ away from $\frac{1}{8}$.

 


	\item (\textit{AHSME 1999}) Find the sum of all prime numbers between $1$ and $100$ that are simultaneously $1$ greater than a multiple of $4$ and $1$ less than a multiple of $5$. 


 $\mathrm{(A) \ } 118 \qquad \mathrm{(B) \ }137 \qquad \mathrm{(C) \ } 158 \qquad \mathrm{(D) \ } 187 \qquad \mathrm{(E) \ } 245$

 \textbf{Answer: }A



\textbf{Solution}

Numbers that are $1$ less than a multiple of $5$ all end in $4$ or $9$.

 No prime number ends in $4$, since all numbers that end in $4$ are divisible by $2$.  Thus, we are only looking for numbers that end in $9$.

 Writing down the ten numbers that so far qualify, we get $9, 19, 29, 39, 49, 59, 69, 79, 89, 99$.

 Crossing off multiples of $3$ gives $19, 29, 49, 59, 79, 89$.

 Crossing off numbers that are not $1$ more than a multiple of $4$ (in other words, numbers that are $1$ less than a multiple of $4$, since all numbers are odd), we get:

 $29, 49, 89.$

 Noting that $49$ is not prime, we have only $29$ and $89$, which give a sum of $118$, so the answer is $\boxed{A}$.

 


	\item (\textit{AHSME 1999}) The marked price of a book was $30 \%$ less than the suggested retail price. Alice purchased the book for half the marked price at a Fiftieth Anniversary sale. What percent of the suggested retail price did Alice pay?


 $\mathrm{(A) \ }25 \% \qquad \mathrm{(B) \ }30 \% \qquad \mathrm{(C) \ }35 \% \qquad \mathrm{(D) \ }60 \% \qquad \mathrm{(E) \ }65 \%$

 \textbf{Answer: }C



\textbf{Solution}

Without loss of generality, let's assume that the retail price was $100$ USD.

 The marked price of the book is $30 \%$ off of $100$ which is equal to $100-100(0.3)=70.$

 Half of that marked price is $0.5(70)=35.$

 Therefore the percent Alice payed of the suggested retail price is $35/100=\boxed{\mathrm{(C) } 35 \%}.$



	\item (\textit{AHSME 1999}) What is the sum of the digits of the decimal form of the product $2^{1999} \cdot 5^{2001}$?


 $\mathrm{(A) \ }2 \qquad \mathrm{(B) \ }4 \qquad \mathrm{(C) \ }5 \qquad \mathrm{(D) \ }7 \qquad \mathrm{(E) \ }10$

 \textbf{Answer: }D



\textbf{Solution}

$2^{1999}\cdot5^{2001}=2^{1999}\cdot5^{1999}\cdot5^{2}=25\cdot10^{1999}$, a number with the digits "25" followed by 1999 zeros. The sum of the digits in the decimal form would be $2+5=7$, thus making the answer $\boxed{\text{D}}$.

 


	\item (\textit{AHSME 1999}) What is the largest number of acute angles that a convex hexagon can have? 


 $\mathrm{(A) \ }2 \qquad \mathrm{(B) \ }3 \qquad \mathrm{(C) \ }4 \qquad \mathrm{(D) \ }5 \qquad \mathrm{(E) \ }6$

 \textbf{Answer: }B



\textbf{Solution}

The sum of the interior angles of a hexagon is $720$ degrees.  In a convex polygon, each angle must be strictly less than $180$ degrees.  

 Six acute angles can only sum to less than $90\cdot 6 = 540$ degrees, so six acute angles could not form a hexagon.

 Five acute angles and one obtuse angle can only sum to less than $90\cdot 5 + 180 = 630$ degrees, so these angles could not form a hexagon.

 Four acute angles and two obtuse angles can only sum to less than $90\cdot 4 + 180\cdot 2 = 720$ degrees.  This is a strict inequality, so these angles could not form a hexagon.  (The limiting figure would be four right angles and two straight angles, which would really be a square with two "extra" points on two sides to form the straight angles.)

 Three acute angles and three obtuse angles work.  For example, if you pick three acute angles of $80$ degrees, the three obtuse angles would be $160$ degrees and give a sum of $80\cdot 3 + 160\cdot 3 = 720$ degrees, which is a genuine hexagon.  Thus, the answer is $\boxed{(B) 3}$



	\item (\textit{AHSME 1999}) At the end of 1994 Walter was half as old as his grandmother. The sum of the years in which they were born is 3838. How old will Walter be at the end of 1999?


 $\mathrm{(A) \ } 48 \qquad \mathrm{(B) \ }49 \qquad \mathrm{(C) \ }53 \qquad \mathrm{(D) \ }55 \qquad \mathrm{(E) \ } 101$

 \textbf{Answer: }D



\textbf{Solution}

In $1994$, if Water is $x$ years old, then Walter's grandmother is $2x$ years old.

 This means that Walter was born in $1994 - x$, and Walter's grandmother was born in $1994 - 2x$.

 The sum of those years is $3838$, so we have:

 $1994 - x + 1994 - 2x = 3838$
$3988 - 3x = 3838$
$x = 50$

 If Walter is $50$ years old in $1994$, then he will be $55$ years old in $1999$, thus giving answer $\boxed{D}$



	\item (\textit{AHSME 1999}) Before Ashley started a three-hour drive, her car's odometer reading was 29792, a palindrome. (A palindrome is a number that reads the same way from left to right as it does from right to left). At her destination, the odometer reading was another palindrome. If Ashley never exceeded the speed limit of 75 miles per hour, which of the following was her greatest possible average speed?


 $\mathrm{(A) \ } 33\frac 13 \qquad \mathrm{(B) \ }53\frac 13 \qquad \mathrm{(C) \ }66\frac 23 \qquad \mathrm{(D) \ }70\frac 13 \qquad \mathrm{(E) \ } 74\frac 13$

 \textbf{Answer: }D



\textbf{Solution}

Ashley could have traveled at most $3\times 75 = 225$ miles.

 Each $5$-digit palindrome is uniquely determined by its first three digits. The next palindromes are 29892, 29992, 30003, and 30103. We may note that $30103-29792 = 311 > 225$, so this number and all larger ones are too large.

 On the other hand, $30003 - 29792 = 211 \leq 225$, thus this is the number on the odometer, and the average speed is $\frac{211}3 = \boxed{(D) 70 \frac 13}$.

 


	\item (\textit{AHSME 1999}) A sealed envelope contains a card with a single digit on it. Three of the following statements are true, and the other is false.


 I. The digit is 1.


 II. the digit is not 2. 


 III. The digit is 3.


 IV. The digit is not 4.


 Which one of the following must necessarily be correct?


 $\mathrm{(A) \ I\ is\ true} \qquad \mathrm{(B) \ I\ is\ false} \qquad \mathrm{(C) \ II\ is\ true} \qquad \mathrm{(D) \ III\ is\ true} \qquad \mathrm{(E) \ IV\ is\ false}$

 \textbf{Answer: }C



\textbf{Solution}

Three of the statements are correct, and only one digit is on the card. Thus, one of I and III are false. Therefore, II and IV must both be true. The answer is therefore $\boxed{\textbf{(C)}}$.

 


	\item (\textit{AHSME 1999}) The student lockers at Olympic High are numbered consecutively beginning with locker number $1$. The plastic digits used to number the lockers cost two cents apiece. Thus, it costs two cents to label locker number $9$ and four cents to label locker number $10$. If it costs $137.94$ to label all the lockers, how many lockers are there at the school?


 $\mathrm{(A) \ }2001 \qquad \mathrm{(B) \ }2010 \qquad \mathrm{(C) \ }2100 \qquad \mathrm{(D) \ }2726 \qquad \mathrm{(E) \ }6897$

 \textbf{Answer: }A



\textbf{Solution 1}

Since all answers are over $2000$, work backwards and find the cost of the first $1999$ lockers.  The first $9$ lockers cost $0.18$ dollars, while the next $90$ lockers cost $0.04\cdot 90 = 3.60$.  Lockers $100$ through $999$ cost $0.06\cdot 900 = 54.00$, and lockers $1000$ through $1999$ inclusive cost $0.08\cdot 1000 = 80.00$.

 This gives a total cost of $0.18 + 3.60 + 54.00 + 80.00 = 137.78$.  There are $137.94 - 137.78 = 0.16$ dollars left over, which is enough for $8$ digits, or $2$ more four digit lockers.  These lockers are $2000$ and $2001$, leading to answer $\boxed{\textbf{(A)}}$.

 


	\item (\textit{AHSME 1999}) What is the maximum number of points of intersection of the graphs of two different fourth degree polynomial functions $y = p(x)$ and $y = q(x)$, each with leading coefficient $1$?


 $\textrm{(A)} \ 1 \qquad \textrm{(B)} \ 2 \qquad \textrm{(C)} \ 3 \qquad \textrm{(D)} \ 4 \qquad \textrm{(E)} \ 8$

 \textbf{Answer: }C



\textbf{Solution}

The intersections of the two polynomials, $p(x)$ and $q(x)$, are precisely the roots of the equation $p(x)=q(x) \rightarrow p(x) - q(x) = 0$. Since the leading coefficients of both polynomials are $1$, the degree of $p(x) - q(x) = 0$ is at most three, and the maximum point of intersection is three, because a third degree polynomial can have at most three roots. Thus, the answer is $\boxed{C}$.

 


	\item (\textit{AHSME 1999}) Define a sequence of real numbers $a_1, a_2, a_3, \ldots$ by $a_1 = 1$ and $a_{n+1}^3 = 99a_n^3$ for all $n \ge 1$. Then $a_{100}$ equals


 $\mathrm{(A) \ } 33^{33} \qquad \mathrm{(B) \ } 33^{99} \qquad \mathrm{(C) \ } 99^{33} \qquad \mathrm{(D) \ }99^{99} \qquad \mathrm{(E) \ none\ of\ the\ above}$

 \textbf{Answer: }C



\textbf{Solution}

We rearrange to get $\dfrac{a_{n+1}}{a_n} = \sqrt[3]{99}$. Thus we get $\dfrac{a_{n+1}}{a_n} = \sqrt[3]{99}$, $\dfrac{a_{n}}{a_{n-1}} = \sqrt[3]{99}$, and so on. Multiplying them all gives $\dfrac{a_{n+1}}{a_1} = (\sqrt[3]{99})^{n}$. Plugging in $n = 99$ and $a_1 = 1$, $a_{100} = (\sqrt[3]{99})^{99} = 99^{33}$, so the answer is $\textbf{(C)}$.

 


	\item (\textit{AHSME 1999}) Four girls - Mary, Aline, Tina, and Hana - sang songs in a concert as trios, with one girl sitting out each time. Hanna sang 7 songs, which was more than any other girl, and Mary sang 4 songs, which was fewer than any other girl. How many songs did these trios sing? 


 $\mathrm{(A) \ 7 } \qquad \mathrm{(B) \ 8 } \qquad \mathrm{(C) \ 9 } \qquad \mathrm{(D) \ 10 } \qquad \mathrm{(E) \ 11 }$

 \textbf{Answer: }A



\textbf{Solution}

Alina and Tina must sing more than $4$, but less than $7$, songs.  Therefore, Alina sang $5$ or $6$ songs, and Tina sang $5$ or $6$ songs, with $4$ possible combinations.

 
However, since every song is a trio, if you add up all the numbers of songs a person sang for all four singers, it must be divisible by $3$.  Thus, $7 + 4 + a + t$ must be divisible by $3$, and $a+t - 1$ must be divisible by $3$.  

 
Since $10 \le a+t \le 12$ by the bounds of $5 \le a, t \le 6$, we have $9 \le a + t - 1 \le 11$.  Because the middle number must be a multiple of $3$, we set $a + t - 1 = 9$, leading to $a + t = 10$ and $a = t = 5$.

 
So, the four girls sang $4, 5, 5$ and $7$ songs in trios.  That means there were $\frac{4 + 5 + 5 + 7}{3} = 7$ songs sung, and the answer is $\boxed{A}$.

 


	\item (\textit{AHSME 1999}) Let $x$ be a real number such that $\sec x - \tan x = 2$. Then $\sec x + \tan x =$
$\mathrm{(A) \ } 0.1 \qquad \mathrm{(B) \ } 0.2 \qquad \mathrm{(C) \ } 0.3 \qquad \mathrm{(D) \ } 0.4 \qquad \mathrm{(E) \ } 0.5$

 \textbf{Answer: }E



\textbf{Solution 1 (Fastest)}

$(\sec x - \tan x)(\sec x + \tan x) = \sec^{2} x - \tan^{2} x = 1$, so $\sec x + \tan x = \boxed{\textbf{(E)}\ 0.5}$.

 



\textbf{Solution 2 (Alternate, Slightly Longer)}

Note that $\sec x - \tan x = (1-\sin x)/\cos x$ and $\sec x + \tan x = (1+\sin x)/\cos x$. Let $(1+\sin x)/\cos x = y$. Multiplying, we get $(1-\sin^{2}x)/\cos^{2}x = 1$.Then, $2y = 1$. $\sec x + \tan x =  \boxed{\textbf{(E)}\ 0.5}$. ~songmath20  Edited 5.1.2023

 


	\item (\textit{AHSME 1999}) What is the radius of a circle inscribed in a rhombus with diagonals of length $10$ and $24$?


 $\mathrm{(A) \ }4 \qquad \mathrm{(B) \ }\frac {58}{13} \qquad \mathrm{(C) \ }\frac{60}{13} \qquad \mathrm{(D) \ }5 \qquad \mathrm{(E) \ }6$

 \textbf{Answer: }C



\textbf{Solution}

Let $d_1=10$ and $d_2=24$ be the lengths of the diagonals, $a$ the side, and $r$ the radius of the inscribed circle.

 Using \href{https://artofproblemsolving.com/wiki/index.php?title=Pythagorean\_theorem}{Pythagorean theorem} we can compute $a=\sqrt{ (d_1/2)^2 + (d_2/2)^2 }=13$.

 We can now express the area of the rhombus in two different ways: as $d_1 d_2 / 2$, and as $2ar$. Solving $d_1 d_2 / 2 = 2ar$ for $r$ we get $r=\boxed{\frac{60}{13}}$.

 (The first formula computes the area as one half of the circumscribed rectangle whose sides are parallel to the diagonals. The second one comes from the fact that we can divide the rhombus into $4$ equal triangles, and in those the height on the side $a$ is equal to $r$. See pictures below.)

 
                    \begin{figure}[H]
                        \centering
                        \includegraphics[width=0.5\textwidth]{Solution/AHSME_I_1999_16_Solution_Figure_1.png}
                    \end{figure}
                    

                    \begin{figure}[H]
                        \centering
                        \includegraphics[width=0.5\textwidth]{Solution/AHSME_I_1999_16_Solution_Figure_2.png}
                    \end{figure}
                    



	\item (\textit{AHSME 1999}) Let $P(x)$ be a polynomial such that when $P(x)$ is divided by $x-19$, the remainder is $99$, and when $P(x)$ is divided by $x - 99$, the remainder is $19$. What is the remainder when $P(x)$ is divided by $(x-19)(x-99)$?


 $\mathrm{(A) \ } -x + 80 \qquad \mathrm{(B) \ } x + 80 \qquad \mathrm{(C) \ } -x + 118 \qquad \mathrm{(D) \ } x + 118 \qquad \mathrm{(E) \ } 0$

 \textbf{Answer: }C



\textbf{Solution 1}

According to the problem statement, there are polynomials $Q(x)$ and $R(x)$ such that $P(x) = Q(x)(x-19) + 99 = R(x)(x-99) + 19$.

 From the last equality we get $Q(x)(x-19) + 80 = R(x)(x-99)$.

 The value $x=99$ is a root of the polynomial on the right hand side, therefore it must be a root of the one on the left hand side as well. Substituting, we get $Q(99)(99-19) + 80 = 0$, from which $Q(99)=-1$. This means that $99$ is a root of the polynomial $Q(x)+1$. In other words, there is a polynomial $S(x)$ such that $Q(x)+1 = S(x)(x-99)$.

 Substituting this into the original formula for $P(x)$ we get 
 \[P(x) = Q(x)(x-19) + 99 = (S(x)(x-99) - 1)(x-19) + 99 =\]
 \[= S(x)(x-99)(x-19) - (x-19) + 99\]

 Therefore when $P(x)$ is divided by $(x-19)(x-99)$, the remainder is $\boxed{-x + 118}$.

 



\textbf{Solution 2}

Since the divisor $(x-19)(x-99)$ is a quadratic, the degree of the remainder is at most linear. We can write $P(x)$ in the form
 \[P(x) = Q(x)(x-19)(x-99) + cx+d\] where $cx+d$ is the remainder.By the Remainder Theorem, plugging in $19$ and $99$ gives us a system of equations.
 \[99c+d = 19\]
 \[19c+d = 99\]

 Solving gives us $c=-1$ and $d = 118$, thus, our answer is $\boxed{ \mathrm{(C) \ }-x+118}$



	\item (\textit{AHSME 1999}) How many zeros does $f(x) = \cos(\log x)$ have on the interval $0 < x < 1$?


 $\mathrm{(A) \ } 0 \qquad \mathrm{(B) \ } 1 \qquad \mathrm{(C) \ } 2 \qquad \mathrm{(D) \ } 10 \qquad \mathrm{(E) \ } \text{infinitely\ many}$

 \textbf{Answer: }E



\textbf{Solution}

For $0 < x < 1$ we have $-\infty < \log x < 0$, and the logarithm is a strictly increasing function on this interval.

 $\cos(t)$ is zero for all $t$ of the form $\frac{\pi}2 + k\pi$, where $k\in\mathbb{Z}$. There are $\boxed{\text{infinitely\ many}}$ such $t$ in $(-\infty,0)$.

 Here's the graph of the function on $(0,1)$:

 
                    \begin{figure}[H]
                        \centering
                        \includegraphics[width=0.5\textwidth]{Solution/AHSME_I_1999_18_Solution_Figure_1.png}
                    \end{figure}
                    

 As we go closer to $0$, the function will more and more wildly oscilate between $-1$ and $1$. This is how it looks like at $(0.0001,0.02)$.

 
                    \begin{figure}[H]
                        \centering
                        \includegraphics[width=0.5\textwidth]{Solution/AHSME_I_1999_18_Solution_Figure_2.png}
                    \end{figure}
                    

 And one more zoom, at $(0.000001,0.0005)$.

 
                    \begin{figure}[H]
                        \centering
                        \includegraphics[width=0.5\textwidth]{Solution/AHSME_I_1999_18_Solution_Figure_3.png}
                    \end{figure}
                    




\textbf{Solution}

If cos(log(x)) = zero, then log(x) = $\pi$/2 + n$\pi$.

 If we consider the limiting case as x approaches zero, log(x) approaches negative infinity. 

 If we consider the other boundary, x equals 1 where log(x) equals zero.

 Due to the intermediate value theorem log(x) must contain all negative real numbers, and thus an infinite number of solutions to log(x) = $\pi$/2 + n$\pi$.This means that cos(x) is zero an infinite number of times giving $\boxed{\text{(E) infinitely\ many}}$.

 -PhysicsMan

 


	\item (\textit{AHSME 1999}) Consider all triangles $ABC$ satisfying in the following conditions: $AB = AC$, $D$ is a point on $\overline{AC}$ for which $\overline{BD} \perp \overline{AC}$, $AC$ and $CD$ are integers, and $BD^{2} = 57$.  Among all such triangles, the smallest possible value of $AC$ is


 
                    \begin{figure}[H]
                        \centering
                        \includegraphics[width=0.5\textwidth]{Problem/AHSME_I_1999_19_Problem_Figure_1.png}
                    \end{figure}
                    
$\textrm{(A)} \ 9 \qquad \textrm{(B)} \ 10 \qquad \textrm{(C)} \ 11 \qquad \textrm{(D)} \ 12 \qquad \textrm{(E)} \ 13$

 \textbf{Answer: }C



\textbf{Solution}

Thus $AD = AC - CD$ and $AB = AC$ are integers. By the \href{https://artofproblemsolving.com/wiki/index.php?title=Pythagorean\_Theorem}{Pythagorean Theorem}, 

 
 \[AD^2 + 57 = AB^2 \Longrightarrow 1 \cdot 57 = 3 \cdot 19 = (AB - AD)(AB + AD).\]

 Thus $AC = AB = \frac {1 + 57}{2} = 29$ or $\frac {3 + 19}{2} = 11 \Longrightarrow \mathrm{(C)}$.

 


	\item (\textit{AHSME 1999}) The sequence $a_{1},a_{2},a_{3},\ldots$ satisfies $a_{1} = 19,a_{9} = 99$, and, for all $n\geq 3$, $a_{n}$ is the arithmetic mean of the first $n - 1$ terms.  Find $a_2$.


 $\textrm{(A)} \ 29 \qquad \textrm{(B)} \ 59 \qquad \textrm{(C)} \ 79 \qquad \textrm{(D)} \ 99 \qquad \textrm{(E)} \ 179$

 \textbf{Answer: }E



\textbf{Solution 1}

Let $m$ be the arithmetic mean of $a_1$ and $a_2$. We can then write $a_1=m-x$ and $a_2=m+x$ for some $x$.

 By definition, $a_3=m$.

 Next, $a_4$ is the mean of $m-x$, $m+x$ and $m$, which is again $m$. 

 Realizing this, one can easily prove by induction that $\forall n\geq 3;~ a_n=m$.

 It follows that $m=a_9=99$. From $19=a_1=m-x$ we get that $x=80$. And thus $a_2 = m+x = \boxed{(E)  179}$.

 



\textbf{Solution 2}

Let $a_1=a$ and $a_2=b$. Then, $a_3=\frac{a+b}{2}$, $a_4=\frac{a+b+\frac{a+b}{2}}{3}=\frac{a+b}{2},$ and so on.

 Since $a_3=a_4$, $a_n=a_3$ for all $n\geq3.$

 Hence, $a_9=\frac{a_1+a_2}{2}=\frac{a+b}{2}=99, a+b=198.$ We also know that $a_1=a=19.$ 

 Subtracting $a_1$ from $198,$ we get $b=a_2=\boxed{(E)  179}.$



	\item (\textit{AHSME 1999}) A circle is circumscribed about a triangle with sides $20,21,$ and $29,$ thus dividing the interior of the circle into four regions. Let $A,B,$ and $C$ be the areas of the non-triangular regions, with $C$ be the largest. Then


 $\mathrm{(A) \ }A+B=C \qquad \mathrm{(B) \ }A+B+210=C \qquad \mathrm{(C) \ }A^2+B^2=C^2 \qquad \mathrm{(D) \ }20A+21B=29C \qquad \mathrm{(E) \ } \frac 1{A^2}+\frac 1{B^2}= \frac 1{C^2}$

 \textbf{Answer: }B



\textbf{Solution}

$20^2 + 21^2 = 841 = 29^2$. Therefore the triangle is a right triangle. But then its hypotenuse is a diameter of the circumcircle, and thus $C$ is exactly one half of the circle. Moreover, the area of the triangle is $\frac{20\cdot 21}{2} = 210$. Therefore the area of the other half of the circumcircle can be expressed as $A+B+210$. Thus the answer is $\boxed{\mathrm{(B)}}$.

 To complete the solution, note that $\mathrm{(A)}$ is clearly false. As $A+B < C$, we have $A^2 + B^2 < (A+B)^2 < C^2$ and thus $\mathrm{(C)}$ is false. Similarly $20A + 21B < 21(A+B) < 21C < 29C$, thus $\mathrm{(D)}$ is false. And finally, since $0<A<C$, $\frac 1{C^2} < \frac1{A^2} < \frac 1{A^2} + \frac 1{B^2}$, thus $\mathrm{(E)}$ is false as well.

 



\textbf{Solution 2 (Alternative to realize that the triangle is Right)}

Copy the link into the browser to see a diagram (I don't know how to link it and I tried using the brackets, but it didn't work)\href{https://geogebra.org/classic/psg3ugzm}{https://geogebra.org/classic/psg3ugzm}

 Let $\circ{O}$ be the circumcircle of $\triangle{XYZ}$ in the problem, and let the circle have a radius of $r$. Let $XY=20, YZ=21, XZ=29$.

 Using the law of cosines: $29^2=20^2+21^2-2*20*21*\cos{XYZ}$.

 Thus $\cos{XYZ}=0$, thus the triangle is right. Thus follow as above: "Moreover, the area of the triangle is $\frac{20\cdot 21}{2} = 210$. Therefore the area of the other half of the circumcircle can be expressed as $A+B+210$. Thus the answer is $\boxed{\mathrm{(B)}}$" (I quoted the solution above to show you where to continue).

 Fixed the link ~\href{https://artofproblemsolving.com/wiki/index.php?title=User:Yiyj1&action=edit&redlink=1}{Yiyj1}

 


	\item (\textit{AHSME 1999}) The graphs of $y = -|x-a| + b$ and $y = |x-c| + d$ intersect at points $(2,5)$ and $(8,3)$. Find $a+c$.


 $\mathrm{(A) \ } 7 \qquad \mathrm{(B) \ } 8 \qquad \mathrm{(C) \ } 10 \qquad \mathrm{(D) \ } 13\qquad \mathrm{(E) \ } 18$

 \textbf{Answer: }C



\textbf{Solution 2}

Plug the coordinates into the two equations and set the equations equal to each other. The $b$ and $d$ terms cancel out and you're left with just the following.

 $|3-a| - |2-a| = |2-c| - |3-c|$

 If we solve individually for the values of $a$ by setting both absolute values on the right to positive and setting one absolute value on the left to a negative, $a = 4$. Doing the same with $c$, we see that $c = 6$. So, $a+c=10$.

 


	\item (\textit{AHSME 1999}) The equiangular convex hexagon $ABCDEF$ has $AB = 1, BC = 4, CD = 2,$ and $DE = 4.$ The area of the hexagon is 


 $\mathrm{(A) \ } \frac {15}2\sqrt{3} \qquad \mathrm{(B) \ }9\sqrt{3} \qquad \mathrm{(C) \ }16 \qquad \mathrm{(D) \ }\frac{39}4\sqrt{3} \qquad \mathrm{(E) \ } \frac{43}4\sqrt{3}$

 \textbf{Answer: }E



\textbf{Solution 1}

Equiangularity means that each internal angle must be exactly $120^\circ$.The information given by the problem statement looks as follows:

 
                    \begin{figure}[H]
                        \centering
                        \includegraphics[width=0.5\textwidth]{Solution/AHSME_I_1999_23_Solution_Figure_1.png}
                    \end{figure}
                    

 We can now place this incomplete polygon onto a triangular grid, finish it, compute its area in unit triangles, and multiply the result by the area of the unit triangle.

 
                    \begin{figure}[H]
                        \centering
                        \includegraphics[width=0.5\textwidth]{Solution/AHSME_I_1999_23_Solution_Figure_2.png}
                    \end{figure}
                    

 We see that the figure contains $43$ unit triangles, and therefore its area is $\boxed{\frac{43\sqrt{3}}4}$.

 



\textbf{Solution 2}


                    \begin{figure}[H]
                        \centering
                        \includegraphics[width=0.5\textwidth]{Solution/AHSME_I_1999_23_Solution_Figure_3.png}
                    \end{figure}
                    

 An equiangular hexagon can be made by drawing an equilateral triangle and cutting out smaller triangles from the corners. Labeling the triangle $X Y$ and $Z$ and drawing $AB$ of length one will remove one equilateral triangle of side length $1$, and drawing $CD$ will take out another equilateral triangle of side length $2$.Labeling the other sides of the smaller equilateral triangles, we can find that $XY$, or the side length of the equilateral triangle is $7$. Now, because we know what the side length of the triangle is, what $DY$ is, and it is given that $DE$ is $4$, we can find the length of $EZ$, $7-4-2=1$. Now, to calculate the area of the hexagon we can simply subtract the area of the smaller equilateral triangles from the larger equilateral triangle. The areas of the smaller equilateral triangles are $\frac{1^2\sqrt{3}}{4}\implies\frac{1\sqrt{3}}{4}$, and $\frac{2^2\sqrt{3}}{4}\implies\frac{4\sqrt{3}}{4}\implies\sqrt{3}$ and the area of the large equilateral triangle is $\frac{7^2\sqrt{3}}{4}\implies\frac{49\sqrt{3}}{4}$ so the area of the hexagon would be $\frac{49\sqrt{3}}{4}-\frac{\sqrt {3}}{4}-\frac{\sqrt{3}}{4}-\sqrt{3}\implies\boxed{\frac{43\sqrt{3}}{4}}$



	\item (\textit{AHSME 1999}) Six points on a circle are given. Four of the chords joining pairs of the six points are selected at random. What is the probability that the four chords form a convex quadrilateral?


 $\mathrm{(A) \ } \frac 1{15} \qquad \mathrm{(B) \ } \frac 1{91} \qquad \mathrm{(C) \ } \frac 1{273} \qquad \mathrm{(D) \ } \frac 1{455} \qquad \mathrm{(E) \ } \frac 1{1365}$

 \textbf{Answer: }B



\textbf{Solution}

There are 15 chords on the circle, which is found by using the diagonal formula n(n-3)/2 where n is 6, and adding 6 adjacent pairs of points. There are 15 choose 4 ways to choose chords, which is 15*14*13*12 divided by 4*3*2*1 = 91 * 15. Choosing four chords that form a convex quadrilateral is the same as choosing four points, since there is only one convex quadrilateral that can be formed with those points. So, the favorable outcomes is 6 Choose 4 which equals 6 Choose 2 which equals 6*5 divided by 2 which equals 15. Therefore, the answer is 15/91*15 = 1/91 = Option B

 


	\item (\textit{AHSME 1999}) There are unique integers $a_{2},a_{3},a_{4},a_{5},a_{6},a_{7}$ such that


 
 \[\frac {5}{7} = \frac {a_{2}}{2!} + \frac {a_{3}}{3!} + \frac {a_{4}}{4!} + \frac {a_{5}}{5!} + \frac {a_{6}}{6!} + \frac {a_{7}}{7!}\]


 where $0\leq a_{i} < i$ for $i = 2,3,\ldots,7$.  Find $a_{2} + a_{3} + a_{4} + a_{5} + a_{6} + a_{7}$.


 $\textrm{(A)} \ 8 \qquad \textrm{(B)} \ 9 \qquad \textrm{(C)} \ 10 \qquad \textrm{(D)} \ 11 \qquad \textrm{(E)} \ 12$

 \textbf{Answer: }B



\textbf{Solution 1(Modular Functions)}

Multiply out the $7!$ to get 

 
 \[5 \cdot 6! = (3 \cdot 4 \cdots 7)a_2 + (4 \cdots 7)a_3 + (5 \cdot 6 \cdot 7)a_4 + 42a_5 + 7a_6 + a_7 .\]

 By \href{https://artofproblemsolving.com/wiki/index.php?title=Wilson\%27s\_Theorem}{Wilson's Theorem} (or by straightforward division), $a_7 + 7(a_6 + 6a_5 + \cdots) \equiv 5 \cdot 6! \equiv -5 \equiv 2 \pmod{7}$, so $a_7 = 2$. Then we move $a_7$ to the left and divide through by $7$ to obtain 

 
 \[\frac{5 \cdot 6!-2}{7} = 514 = 360a_2 + 120a_3 + 30a_4 + 6a_5 + a_6.\] 

 We then repeat this procedure $\pmod{6}$, from which it follows that $a_6 \equiv 514 \equiv 4 \pmod{6}$, and so forth. Continuing, we find the unique solution to be $(a_2, a_3, a_4, a_5, a_6, a_7) = (1,1,1,0,4,2)$ (uniqueness is assured by the \href{https://artofproblemsolving.com/wiki/index.php?title=Division\_Theorem}{Division Theorem}). The answer is $9 \Longrightarrow \boxed{\mathrm{(B)}}$.

 



\textbf{Solution 2(Basic Algebra and Bashing)}

We start by multiplying both sides by $7!$, and we get:
 \[3600=2520a_2+840a_3+210a_4+42a_5+7a_6+a_7\]After doing some guess and check, we find that the answer is $\boxed{\textbf{(B)~9}}$.

 



\textbf{Solution 3 (The Easiest and Most Intuitive)}

Let's clear up the fractions: 
 \[\frac{5}{7}=\frac{2520a_2+840a_3+210a_4+42a_5+7a_6+a_7}{7!}\]
 \[3600=2520a_2+840a_3+210a_4+42a_5+7a_6+a_7\]Notice that if we divide everything by $7$ then we would have:
 \[514+\frac{2}{7}=360a_2+120a_3+30a_4+6a_5+a_6+\frac{1}{7}a_7\]Since $0 \le a_7<7$ and $a_7$ must be an integer, then we have $\frac{2}{7}=\frac{1}{7}a_7$, so $a_7=2$.

 Similarly, if we divide everything by $6$, then we would have: 
 \[85+\frac{4}{6}=60a_2+20a_3+5a_4+a_5+\frac{1}{6}a_6\]Again, since $0 \le a_6<6$ and $a_6$ must be an integer, we have $\frac{4}{6}=\frac{1}{6}a_6$, so $a_6=4$.

 The pattern repeats itself, so in the end we have $a_2=1$, $a_3=1$, $a_4=1$, $a_5=0$, $a_6=4$, $a_7=2$. So $a_2+a_3+a_4+a_5+a_6+a_7=\boxed{\textbf{(B)~9}}$




\textbf{Solution 4}

By multiplying both sides by $7$ we get 

 
 \[5 = \frac72 a_2 + \frac76 a_3+ \frac{7}{24} a_4 + \frac{7}{120} a_5 + \frac{7}{720} a_6 + \frac{a_7}{720}\]

 since $0 \le a_2 < 2$, if $a_2 = 0$ the rest of the right hand side will not add up to be $5$, so $a_2 = 1$

 \[\frac32 = \frac76 a_3+ \frac{7}{24} a_4 + \frac{7}{120} a_5 + \frac{7}{720} a_6 + \frac{a_7}{720}\]

 If $a_3 = 2$, $\frac76 a_3 = \frac73 > \frac32$, so $a_3 = 1$

 \[\frac13 = \frac{7}{24} a_4 + \frac{7}{120} a_5 + \frac{7}{720} a_6 + \frac{a_7}{720}\]

 If $a_4 = 2$, $\frac{7}{24} a_4 = \frac{7}{12} > \frac13$, so $a_4 = 1$

 \[\frac{1}{24} = \frac{7}{120} a_5 + \frac{7}{720} a_6 + \frac{a_7}{720}\]

 If $a_5 = 1$, $\frac{7}{120} a_5 = \frac{7}{120} > \frac{1}{24}$, so $a_5 = 0$

 \[\frac{1}{24} =  \frac{7}{720} a_6 + \frac{a_7}{720}\]
$30 = 7 a_6 + a_7$. Since $a_7 < 7$, $a_7 = 2$ and $a_6=4$

 Therefore, $a_{2} + a_{3} + a_{4} + a_{5} + a_{6} + a_{7} = 1 + 1 + 1 + 0 + 4 + 2= \boxed{\textbf{(B) } 9}$.

 


	\item (\textit{AHSME 1999}) Three non-overlapping regular plane polygons, at least two of which are congruent, all have sides of length $1$. The polygons meet at a point $A$ in such a way that the sum of the three interior angles at $A$ is $360^{\circ}$. Thus the three polygons form a new polygon with $A$ as an interior point. What is the largest possible perimeter that this polygon can have? 


 $\mathrm{(A) \ }12 \qquad \mathrm{(B) \ }14 \qquad \mathrm{(C) \ }18 \qquad \mathrm{(D) \ }21 \qquad \mathrm{(E) \ } 24$

 \textbf{Answer: }D



\textbf{Solution}

We are looking for three regular polygons such that the sum of their internal angle sizes is exactly $360^{\circ}$. 

 Let the number of sides in our polygons be $3\leq a,b,c$. From each of the polygons, two sides touch the other two, and the remaining sides are on the perimeter. Therefore the answer to our problem is the value $(a-2)+(b-2)+(c-2) = (a+b+c)-6$.

 The integral angle of a regular $k$-gon is $180 \frac{k-2}k$. Therefore we are looking for integer solutions to:

 
 \[360 = 180\left( \frac{a-2}a + \frac{b-2}b + \frac{c-2}c \right)\]

 Which can be simplified to:

 
 \[2 = \left( \frac{a-2}a + \frac{b-2}b + \frac{c-2}c \right)\]

 Furthermore, we know that two of the polygons are congruent, thus WLOG $a=c$. Our equation now becomes

 
 \[2 = \left( 2\cdot\frac{a-2}a + \frac{b-2}b \right)\]

 Multiply both sides by $ab$ and simplify to get $ab - 4b - 2a = 0$.

 Using the standard technique for \href{https://artofproblemsolving.com/wiki/index.php?title=Diophantine\_equation}{Diophantine equations}, we can add $8$ to both sides and rewrite the equation as $(a-4)(b-2)=8$.

 Remembering that $a,b\geq 3$ the only valid options for $(a-4,b-2)$ are: $(1,8)$, $(2,4)$, $(4,2)$, and $(8,1)$.

 These correspond to the following pairs $(a,b)$: $(5,10)$, $(6,6)$, $(8,4)$, and $(12,3)$. 

 The perimeters of the resulting polygon for these four cases are $14$, $12$, $14$, and $\boxed{21}$.

 



\textbf{Solution 2}

We want to maximize the number of sides of the two congruent polygons, so we need to make the third polygon have the fewest number of sides possible, i.e. a triangle. The interior angle measure of the two congruent polygons is therefore $\frac{360-60}{2}=150$ degrees, so they are dodecagons. Of all the $12 + 12 + 3 = 27$ sides, six of them are not part of the perimeter of the resulting polygon, so the resulting polygon has $\boxed{21}$ sides.

 


	\item (\textit{AHSME 1999}) In triangle $ABC$, $3 \sin A + 4 \cos B = 6$ and $4 \sin B + 3 \cos A = 1$. Then $\angle C$ in degrees is 


 $\mathrm{(A) \ }30 \qquad \mathrm{(B) \ }60 \qquad \mathrm{(C) \ }90 \qquad \mathrm{(D) \ }120 \qquad \mathrm{(E) \ }150$

 \textbf{Answer: }A



\textbf{Solution}

Square the given equations and add (simplifying with the Pythagorean identity $\sin^2 x + \cos^2 x = 1$):

 
 \begin{align*} 9\sin^2 A + 16\cos^2 B + 24 \sin A \cos B & = 36 \\ + 9\cos^2 A + 16\sin^2 B + 24 \sin B \cos A & = 1 \\ \Longrightarrow 25 + 24(\sin A \cos B + \sin B \cos A ) & = 37 \end{align*}


 Thus $\frac 12 = \sin A \cos B + \sin B \cos A$. This is the sine addition identity, so $\frac 12 = \sin (A + B) = \sin (180 - C) = \sin C$. Thus either $C = 30^{\circ}, 150^{\circ}$. 

 If $C = 150$, then $A + B = 30 \Longrightarrow A,B < 30$, and $\sin A < \frac 12, \cos A < 1$. The first equation implies $6 = 3 \sin A + 4\cos B < 3\left(\frac 12\right) + 4(1) = 5.5 < 6$, which is a contradiction; thus $C = 30 \Longrightarrow \mathrm{(A)}$.

 \textbf{Supplement} Adding the two given equations gives $3(\sin A + \cos A)+4(\sin B + \cos B) =7$

 If $A+B = 30^{\circ}$, then $A$ and $B$ are both acute. When $A$ and $B$ are both acute, $\sin A + \cos A > 1$ and $\sin B + \cos B>1$. Then $3(\sin A + \cos A)+4(\sin B + \cos B) >7$. This contradicts the equation $3(\sin A + \cos A)+4(\sin B + \cos B) =7$. Therefore, $A+B \neq 30^{\circ}$, so $A+B = 150^{\circ}$ and $C = 30^{\circ}$.

 


	\item (\textit{AHSME 1999}) Let $x_1, x_2, \ldots , x_n$ be a sequence of integers such that


 $\text{(i)}$ $-1 \le x_i \le 2$ $\text{for}$ $i = 1,2, \ldots n$
$\text{(ii)}$ $x_1 + \cdots + x_n = 19$; $\text{and}$
$\text{(iii)}$ $x_1^2 + x_2^2 + \cdots + x_n^2 = 99$.


 Let $m$ and $M$ be the minimal and maximal possible values of $x_1^3 + \cdots + x_n^3$, respectively. Then $\frac Mm =$
$\mathrm{(A) \ }3 \qquad \mathrm{(B) \ }4 \qquad \mathrm{(C) \ }5 \qquad \mathrm{(D) \ }6 \qquad \mathrm{(E) \ }7$

 \textbf{Answer: }E



\textbf{Solution 1}

Clearly, we can ignore the possibility that some $x_i$ are zero, as adding/removing such variables does not change the truth value of any condition, nor does it change the value of the sum of cubes. Thus we'll only consider $x_i\in\{-1,1,2\}$. 

 Also, order of the $x_i$ does not matter, so we are only interested in the counts of the variables of each type. Let $a$ of the $x_i$ be equal to $-1$, $b$ equal to $1$, and $c$ equal to $2$.

 The conditions (ii) and (iii) simplify to: 
(ii) $2c + b - a = 19$ 
(iii) $4c + a + b = 99$ 
and we want to find the maximum and minimum of $2^3c + 1^3b + (-1)^3a = 8c + b - a$ over all non-negative solutions of the above two equations.

 Subtracting twice (ii) from (iii) we get $b = 3a-61$. By entering that into one of the two equations and simplifying we get $c=40-a$.

 Thus all the solutions of our system of equations have the form $(a,3a-61,40-a)$.

 As all three variables must be non-negative integers, we have $3a-61 \geq 0 \Rightarrow a \geq 21$ and $40-a\geq 0 \Rightarrow a \leq 40$.

 For $(a,b,c)$ of the form $(a,3a-61,40-a)$ the expression we are maximizing/minimizing simplifies to $320-8a+3a-61-a = 259 - 6a$. Clearly, the maximum is achieved for $a=21$ and the minimum for $a=40$. Their values are $M=133$ and $m=19$, and therefore $\frac Mm = \boxed{7}$.

 \textbf{Note} The minimum is achieved for 
 \[(\underbrace{-1,\dots,-1}_{40},\underbrace{1,\dots,1}_{59})\]The maximum is achieved for
 \[(\underbrace{-1,\dots,-1}_{21},1,1,\underbrace{2,\dots,2}_{19})\]




\textbf{Solution 2}

As said in Solution 1, $x_i=0$ can be ignored, and only $x_i = -1, 1, 2$ need to be considered.

 \textbf{Minimum} To minimize $x_1^3 + \cdots + x_n^3$, there are no $2$s and maximize the number of $-1$s. 

 Therefore, the number of $1$s are $19 + \frac{99-19}{2} = 59$, the number of $-1$s are $\frac{99-19}{2} = 40$.

 $x_1^3 + \cdots + x_n^3 = 59 - 40 = 19$

 \textbf{Maximum} Let $a =$ number of $2$s, $b =$ number of $1$s, $c =$ number of $-1$s

 
 \[4a+b+c=99\]
 \[2a+b-c=19\]

 Multiplying the second equation by $2$ gives $4a+2b-2c=38$. By subtracting this equation from the first equation we get $3c-b=61$, $3c=61+b$. As we need to minimize the value of $c$, $c = 21$, $b = 2$, $a = 19$
$x_1^3 + \cdots + x_n^3 = 8 \cdot 19 +2 - 21 = 133$

 Therefore, $\frac Mm = \frac{133}{19}=\boxed{\textbf{(E) } 7}$.

 


	\item (\textit{AHSME 1999}) A tetrahedron with four equilateral triangular faces has a sphere inscribed within it and a sphere circumscribed about it. For each of the four faces, there is a sphere tangent externally to the face at its center and to the circumscribed sphere. A point $P$ is selected at random inside the circumscribed sphere. The probability that $P$ lies inside one of the five small spheres is closest to 


 $\mathrm{(A) \ }0 \qquad \mathrm{(B) \ }0.1 \qquad \mathrm{(C) \ }0.2 \qquad \mathrm{(D) \ }0.3 \qquad \mathrm{(E) \ }0.4$

 \textbf{Answer: }C



\textbf{Solution 1}

Let the radius of the large sphere be $R$, and of the inner sphere $r$. Label the vertices of the tetrahedron $ABCD$, and let $O$ be the center. Then pyramid $[OABC] + [OABD] + [OACD] + [OBCD] = [ABCD]$, where $[\ldots]$ denotes volume; thus $[OABC] = \frac{[ABCD]}{4}$. Since $OABC$ and $ABCD$ are both pyramids that share a common face $ABC$, the ratio of their volumes is the ratio of their altitudes to face $ABC$, so $r = \frac {h_{ABCD}}4$. However, $h_{ABCD} = r + R$, so it follows that $r = \frac {R}{3}$. Then the radius of an external sphere is $\frac{R-r}2 = \frac {R}{3} = r$.

 Since the five described spheres are non-intersecting, it follows that the ratio of the volumes of the spheres is $5 \cdot \left( \frac 13 \right)^3 = \frac{5}{27} \approx 0.2 \Longrightarrow \fbox{C}$.

 



\textbf{Solution 2}

Let the side length of the tetrahedron be $1$.

 The altitude of the equilateral triangle base is $\frac{\sqrt{3}}{2}$. Thus, the distance from the center of the equilateral triangle to its vertex is $\frac{\sqrt{3}}{2} \cdot \frac 23 = \frac{\sqrt{3}}{3}$. Therefore, the altitude of the tetrahedron is $\sqrt{1^2-(\frac{\sqrt{3}}{3})^2} = \frac{\sqrt{6}}{3}$.

 Let the radius of the large sphere be $R$.
 \[(\frac{\sqrt{3}}{3})^2 + (\frac{\sqrt{6}}{3} - R)^2 = R^2\]
 \[R = \frac{\sqrt{6}}{4}\]

 The distance from the center of the tetrahedron to the center of one of the bases is $\frac{\sqrt{6}}{3} - R = \frac{\sqrt{6}}{3} - \frac{\sqrt{6}}{4} = \frac{\sqrt{6}}{12}$

 The distance from the center of the tetrahedron to the center of one of the bases is also the same as the radius of the small sphere in the center of the tetrahedron.

 The radius of the smaller spheres tangent to the large sphere and the tetrahedron is $\frac{\frac{\sqrt{6}}{4}-\frac{\sqrt{6}}{12}}{2} = \frac{\sqrt{6}}{12}$

 Therefore, the probability that $P$ lies inside one of the five small spheres is $\frac{5 \cdot (\frac{\sqrt{6}}{12})^3 \cdot \frac43 \cdot \pi }{(\frac{\sqrt{6}}{4})^3 \cdot \frac43 \cdot \pi} = \frac{5}{27} \approx \boxed{\textbf{(C) } 0.2}$



	\item (\textit{AHSME 1999}) The number of ordered pairs of integers $(m,n)$ for which $mn \ge 0$ and


 
 \[m^3 + n^3 + 99mn = 33^3\]


 is equal to


 $\mathrm{(A) \ }2 \qquad \mathrm{(B) \ } 3\qquad \mathrm{(C) \ } 33\qquad \mathrm{(D) \ }35 \qquad \mathrm{(E) \ } 99$

 \textbf{Answer: }D



\textbf{Solution 1 (symmetric sums)}

We recall the factorization (see \href{https://artofproblemsolving.com/wiki/index.php?title=Elementary\_symmetric\_sum}{elementary symmetric sums})

 
 \[x^3 + y^3 + z^3 - 3xyz = \frac12\cdot(x + y + z)\cdot((x - y)^2 + (y - z)^2 + (z - x)^2)\]

 Setting $x = m,y = n,z = - 33$, we have that either $m + n - 33 = 0$ or $m = n = - 33$ (by the \href{https://artofproblemsolving.com/wiki/index.php?title=Trivial\_Inequality}{Trivial Inequality}). Thus, there are $35 \Longrightarrow \mathrm{(D)}$ solutions satisfying $mn \ge 0$.

 


 



\textbf{Solution 2 (Binomial theorem / symmetric sums)}

Since we see $^3$ and multiples of $3$, let's try the binomial theorem:

 $(m+n)^3 = m^3 + n^3 + 3(m+n)mn$.

 This is almost exactly the main premise of the problem, so next try $m+n=33$:

 $33^3 = m^3 + n^3 + 99mn$

 Since we require $mn \geq 0$,  we have solutions for all $0 \leq m \leq 33$, which has $\mathbf{ 34 }$ solutions.

 If $\max(m, n) \geq 33$, the left side must be greater than the right side, so these are the only possible positive solutions.

 Next, we consider negative solutions. (Also, the multiple choice options confirm that we are missing a solution. Further, since the remaining options are all odd, we get a hint that, due to symmetry, there must be an odd number of solutions with $m=n$.).

 Consider the case $m=n$: $2m^3 + 99m^2 = 33^3$, equivalent to $m^2(2m+99) = 33^3$. If you don't immediately guess $m=-33$ (which works), observe that $m \equiv 0 \pmod {11}$ and also $\pmod 3$.Arithmetic shows that $(-33,-33)$ is a solution for this case, and no other values of $m$ give $m^2(2m+99)  \geq 0$. So we have $\mathbf{1}$ additional solution.

 This gives the correct answer of $34+1=\boxed{\textbf{(D) } 35}$, but we haven't yet proved there are no other solutions..

 Finally, we consider $m \neq n$, when both are negative. This forces$m^3 \equiv -(n^3)$ in both $\pmod {11}$ and $\pmod 3$,which implies (by computing all the modular cubes) that $m \equiv -n$.This further implies that $m \equiv -n \pmod {33}$, due to the modular (Chinese) Remainder Theorem.

 But as $|m-n|$ increases $m^3+n^3$ grows more negative while $99mn$ shrinks toward 0. Thus the sum cannot remain constant, unless one of $m$ or $n$ jumps by 33.

 There are a few remaining possibilities to check before $m^3$ grows far too large, but at this point we can make a nearly confident guess that $\textbf{(E) } 99$ is impossible. :-/

 



\textbf{Solution 3 (Factoring from first principles)}

The coefficients $33$ and $99$ are related, so substitute $p = 33$; hence $m^3 + n^3 + 3mnp = p^3$ or $m^3 + n^3 - p^3 + 3mnp = 0$.Rewrite $m^3 + n^3 = (m+n)^3 - 3m^2n - 3mn^2 = (m+n)^3 - 3mn(m+n)$. The motivation is to combine the $m^3$ and $n^3$ terms into a single cubed term that can be further combined with the $-p^3$.Then \begin{align*}(m+n)^3 - p^3 - 3mn(m+n) + 3mnp &= 0 \\(m+n-p)((m+n)^2 + (m+n)p + p^2) + 3mn(p-m-n) &= 0 \\(m+n-p)((m+n)^2 + (m+n)p + p^2 - 3mn) &= 0 \\(m+n-p)(m^2 + n^2 + p^2 + mp + np - mn) &= 0 \\(m+n-p)*\frac{1}{2} ((m+p)^2 + (n+p)^2 + (m-n)^2) &= 0\end{align*}


 It follows that either $m + n = p$ or $m = n = -p$, and as in Solution 1, there are $\boxed{\textbf{(D) } 35}$ pairs.

 


\end{enumerate}